% ============================================================
%  RAPPORT DE PROJET FINAL — Travel Maroc Agency
%  JOUMAIL Abderrahmane — CDL 1ère année — EST Dakhla 2025-2026
%  Compiler avec : pdflatex (2 passes) ou Overleaf
% ============================================================
\documentclass[12pt,a4paper]{report}

% ── Encodage & langue ──────────────────────────────────────
\usepackage[T1]{fontenc}
\usepackage[utf8]{inputenc}
\usepackage[french]{babel}
\usepackage{lmodern}
\usepackage{microtype}

% ── Géométrie ─────────────────────────────────────────────
\usepackage[left=3cm,right=2.5cm,top=3cm,bottom=2.5cm,headheight=15pt]{geometry}

% ── Couleurs ──────────────────────────────────────────────
\usepackage{xcolor}
\definecolor{estblue}{RGB}{26,35,126}
\definecolor{estgold}{RGB}{255,179,0}
\definecolor{lightblue}{RGB}{232,234,246}
\definecolor{lightgold}{RGB}{255,249,210}
\definecolor{codebg}{RGB}{248,249,250}
\definecolor{codegreen}{RGB}{39,128,39}
\definecolor{codegray}{RGB}{120,120,120}
\definecolor{codered}{RGB}{180,0,0}
\definecolor{darkbg}{RGB}{30,30,30}
\definecolor{tmablue}{RGB}{13,43,85}
\definecolor{tmaorange}{RGB}{249,115,22}

% ── Graphiques & TikZ ─────────────────────────────────────
\usepackage{graphicx}
\graphicspath{{screenshots/}}
\usepackage{tikz}
\usetikzlibrary{arrows.meta,positioning,calc,shapes.geometric,fit,backgrounds,shadows.blur}

% ── Tableaux ──────────────────────────────────────────────
\usepackage{booktabs}
\usepackage{tabularx}
\usepackage{longtable}
\usepackage{array}
\usepackage{multirow}
\renewcommand{\arraystretch}{1.35}

% ── Code & tcolorbox ──────────────────────────────────────
\usepackage{listings}
\usepackage{tcolorbox}
\tcbuselibrary{listings,breakable,skins}

\lstdefinestyle{codestyle}{
  basicstyle=\footnotesize\ttfamily,
  breaklines=true,
  breakatwhitespace=false,
  commentstyle=\color{codegreen}\itshape,
  keywordstyle=\color{estblue}\bfseries,
  numberstyle=\tiny\color{codegray},
  stringstyle=\color{codered},
  numbers=left,
  numbersep=8pt,
  showstringspaces=false,
  tabsize=2,
  frame=none,
  backgroundcolor=\color{codebg},
}

\lstdefinelanguage{PHP}{
  keywords={php,class,public,private,protected,static,function,return,if,else,
    foreach,while,for,echo,new,null,true,false,global,use,namespace,extends,
    implements,abstract,interface,const,self,parent,array,string,int,bool,void},
  morecomment=[l]{//},
  morecomment=[s]{/*}{*/},
  morestring=[b]',
  morestring=[b]",
  sensitive=true
}

% Boîte code PHP (bleu)
\newtcblisting{phpbox}{
  listing only,
  listing engine=listings,
  listing options={style=codestyle,language=PHP},
  colback=codebg, colframe=estblue,
  arc=4pt, boxrule=0.6pt,
  left=4pt, right=4pt, top=3pt, bottom=3pt,
  enhanced, breakable=false,
  before skip=10pt, after skip=10pt,
  before upper={\needspace{12\baselineskip}},
}

% Boîte shell (fond sombre)
\newtcblisting{shellbox}{
  listing only,
  listing engine=listings,
  listing options={style=codestyle,
    basicstyle=\footnotesize\ttfamily\color{white!90},
    backgroundcolor=\color{darkbg},
    commentstyle=\color{green!60!white}\itshape,
    keywordstyle=\color{cyan!70!white}\bfseries,
    stringstyle=\color{orange!80!white}},
  colback=darkbg, colframe=black!70,
  arc=4pt, boxrule=0.6pt,
  left=4pt, right=4pt, top=3pt, bottom=3pt,
  enhanced, breakable=false,
  before skip=10pt, after skip=10pt,
  before upper={\needspace{12\baselineskip}},
}

% Boîte JSON (doré)
\newtcblisting{jsonbox}{
  listing only,
  listing engine=listings,
  listing options={style=codestyle},
  colback=codebg, colframe=estgold,
  arc=4pt, boxrule=0.6pt,
  left=4pt, right=4pt, top=3pt, bottom=3pt,
  enhanced, breakable=false,
  before skip=10pt, after skip=10pt,
  before upper={\needspace{12\baselineskip}},
}

% ── En-têtes & pieds de page ──────────────────────────────
\usepackage{fancyhdr}
\pagestyle{fancy}
\fancyhf{}
\fancyhead[L]{\small\color{estblue}\textit{\leftmark}}
\fancyhead[R]{\small\color{estblue}\textbf{Travel Maroc Agency}}
\fancyfoot[C]{\small\color{estblue}\thepage}
\fancyfoot[R]{\tiny\color{gray}JOUMAIL Abderrahmane --- CDL 2025--2026}
\renewcommand{\headrulewidth}{0.5pt}
\renewcommand{\footrulewidth}{0.3pt}

% ── Formatage chapitres/sections ──────────────────────────
\usepackage{titlesec}

\titleformat{\chapter}[display]
  {\normalfont}
  {\colorbox{estblue}{\parbox{\dimexpr\linewidth-2\fboxsep\relax}{%
      \centering\vspace{5pt}\color{white}\Large\bfseries\chaptertitlename\ \thechapter\vspace{5pt}}}}
  {0pt}
  {\vspace{10pt}\huge\bfseries\color{estblue}\centering}
  [\vspace{6pt}\color{estgold}\rule{\textwidth}{2.2pt}\vspace{2pt}]

\newcommand{\diagramtitle}[1]{%
  \par\medskip
  {\centering\large\bfseries\color{estblue}#1\par}
  \vspace{0.35em}
}

\tikzset{
  diagram frame/.style={rounded corners=10pt, draw=estblue!35, line width=0.9pt,
    fill=white, blur shadow={shadow xshift=1.5pt, shadow yshift=-1.5pt,
    shadow blur steps=5, shadow opacity=20}},
  diagram box/.style={diagram frame, fill=lightblue!35, minimum height=1.65cm,
    text width=4.6cm, align=center, inner sep=10pt, font=\small},
  diagram box gold/.style={diagram frame, fill=lightgold!80!white, draw=estgold!80!black,
    minimum height=1.65cm, text width=4.6cm, align=center, inner sep=10pt, font=\small},
  diagram actor/.style={diagram frame, fill=lightgold, draw=estgold!80!black,
    minimum width=3.2cm, minimum height=1.1cm, align=center, font=\small\bfseries},
  diagram db/.style={diagram frame, fill=estblue!10, minimum width=4.2cm,
    minimum height=1.5cm, align=center, inner sep=9pt, font=\small},
  diagram arrow/.style={-{Stealth[length=3.6mm,width=2.4mm]}, line width=1.4pt, estblue},
  diagram note/.style={font=\scriptsize\itshape, text=gray!80, fill=white,
    inner sep=2pt, align=center},
  diagram small/.style={font=\scriptsize, text=gray!85}
}

\titleformat{\section}
  {\normalfont\Large\bfseries\color{estblue}}
  {\thesection}{0.8em}{}
  [\vspace{-4pt}\color{estblue!35}\rule{\linewidth}{0.5pt}]

\titleformat{\subsection}
  {\normalfont\large\bfseries\color{estblue!80}}
  {\thesubsection}{0.8em}{}

\titleformat{\subsubsection}
  {\normalfont\normalsize\bfseries\color{estblue!65}}
  {\thesubsubsection}{0.8em}{}

\titlespacing{\chapter}{0pt}{0pt}{18pt}
\titlespacing{\section}{0pt}{14pt}{6pt}
\titlespacing{\subsection}{0pt}{10pt}{4pt}

% ── Table des matières ────────────────────────────────────
\usepackage{tocloft}
\renewcommand{\cftchapfont}{\bfseries\color{estblue}}
\renewcommand{\cftchappagefont}{\bfseries\color{estblue}}
\renewcommand{\cftsecfont}{\color{estblue!80}}
\renewcommand{\cftsecpagefont}{\color{estblue!80}}
\renewcommand{\cftchapleader}{\color{estgold}\cftdotfill{\cftchapdotsep}}

% ── Boîtes d'information ──────────────────────────────────
\usepackage{mdframed}

\newmdenv[linecolor=estgold,backgroundcolor=lightgold,
  linewidth=2.5pt,leftline=true,rightline=false,
  topline=false,bottomline=false,
  innerleftmargin=10pt,innertopmargin=6pt,innerbottommargin=6pt]{notebox}

\newmdenv[linecolor=estblue,backgroundcolor=lightblue!50,
  linewidth=2.5pt,leftline=true,rightline=false,
  topline=false,bottomline=false,
  innerleftmargin=10pt,innertopmargin=6pt,innerbottommargin=6pt]{infobox}

% ── Listes ────────────────────────────────────────────────
\usepackage{enumitem}
\setlist[itemize]{leftmargin=*,label={\color{estgold}\textbullet}}
\setlist[enumerate]{leftmargin=*}

% ── Divers ────────────────────────────────────────────────
\usepackage{setspace}
\onehalfspacing
\usepackage{float}
\usepackage{caption}
\usepackage[colorlinks=true,linkcolor=estblue,urlcolor=estblue,citecolor=estblue]{hyperref}
\usepackage{amssymb}
\usepackage{needspace}
\usepackage{etoolbox}

% ── Prévention titres orphelins ───────────────────────────
\raggedbottom
\widowpenalty=10000
\clubpenalty=10000
\displaywidowpenalty=10000
\preto{\chapter}{\needspace{10\baselineskip}}
\preto{\section}{\needspace{6\baselineskip}}
\preto{\subsection}{\needspace{4\baselineskip}}

% =============================================================
\begin{document}
% =============================================================

% ─────────────────────────────────────────────────────────────
%  PAGE DE GARDE
% ─────────────────────────────────────────────────────────────
\begin{titlepage}
\begin{tikzpicture}[remember picture,overlay]
  \fill[estblue] (current page.north west)
    rectangle ([yshift=-4cm] current page.north east);
  \fill[estgold] ([yshift=-4cm] current page.north west)
    rectangle ([yshift=-4.35cm] current page.north east);
  \fill[estblue] (current page.south west)
    rectangle ([yshift=2.2cm] current page.south east);
  \fill[estgold] ([yshift=2.2cm] current page.south west)
    rectangle ([yshift=2.55cm] current page.south east);
  \node[anchor=west]
    at ([xshift=2.3cm,yshift=-2.15cm] current page.north west)
    {\includegraphics[width=5.1cm]{logo_ESTD.jpeg}};
  \node[estgold,font=\normalsize\bfseries,anchor=west]
    at ([xshift=8.0cm,yshift=-2.2cm] current page.north west)
    {Fili\`ere Conception et D\'eveloppement de Logiciel (CDL)};
  \node[white,font=\small,anchor=center]
    at ([yshift=1.35cm] current page.south)
    {Ann\'ee Universitaire 2025--2026 \quad\textbullet\quad
     Module : D\'eveloppement Web Avanc\'e};
\end{tikzpicture}

\vspace{4.5cm}
\begin{center}

\begin{tikzpicture}[scale=1.1]
  \draw[estblue,line width=2.5pt] (0,0) circle (1.4cm);
  \draw[estgold,line width=1.8pt]
    (-0.80,0.025) -- (-0.15,0.725) -- (0.80,0.725) --
    (0.80,-0.725) -- (-0.80,-0.725) -- cycle;
  \draw[estgold,line width=1.2pt] (-0.60,0.255) -- (0.60,0.255);
  \draw[estgold,line width=1.2pt] (-0.60,-0.105) -- (0.60,-0.105);
  \draw[estgold,line width=1.2pt] (-0.60,-0.445) -- (0.07,-0.445);
  \node[estblue,font=\LARGE\bfseries] at (0,-1.8) {Travel};
  \node[estgold,font=\large]          at (0,-2.3) {Maroc Agency};
\end{tikzpicture}

\vspace{1.2cm}
{\color{estblue!40}\rule{0.45\textwidth}{0.6pt}}
\vspace{0.7cm}

{\fontsize{28}{40}\selectfont\bfseries\color{estblue}
RAPPORT DE PROJET\\[4pt]
FINAL}

\vspace{0.45cm}
{\Large\color{estblue!70} Application Web Full-Stack}

\vspace{0.3cm}
{\fontsize{20}{26}\selectfont\bfseries\color{estgold} Travel Maroc Agency}

\vspace{0.3cm}
{\large\itshape\color{gray}
Plateforme de R\'eservation de Voyages et de Gestion d'Agence Touristique}

\vspace{0.9cm}
{\color{estblue!40}\rule{0.45\textwidth}{0.6pt}}
\vspace{1cm}

\renewcommand{\arraystretch}{1.6}
\begin{tabular}{r@{\hspace{8pt}:\hspace{8pt}}l}
  \color{estblue}\textbf{R\'ealis\'e par}  &
    \textbf{JOUMAIL Abderrahmane} \\
  \color{estblue}\textbf{Fili\`ere}        &
    Conception et D\'eveloppement de Logiciel (CDL) \\
  \color{estblue}\textbf{Niveau}           &
    1\textsuperscript{\`ere} ann\'ee \\
  \color{estblue}\textbf{Encadrant}        &
    Mme Ouahi Mariame \\
  \color{estblue}\textbf{Module}           &
    D\'eveloppement Web Avanc\'e \\
  \color{estblue}\textbf{Date de rendu}    &
    Mai 2026 \\
  \color{estblue}\textbf{URL locale}       &
    \texttt{http://localhost/boutique/} \\
\end{tabular}
\renewcommand{\arraystretch}{1.35}

\end{center}
\end{titlepage}

% ─────────────────────────────────────────────────────────────
%  PAGES PRÉLIMINAIRES
% ─────────────────────────────────────────────────────────────
\pagenumbering{roman}

\chapter*{R\'esum\'e}
\addcontentsline{toc}{chapter}{R\'esum\'e}
\thispagestyle{fancy}

Ce rapport pr\'esente la conception et le d\'eveloppement de \textbf{Travel Maroc Agency},
une plateforme web complète de r\'eservation de voyages destin\'ee \`a une agence
touristique bas\'ee \`a Dakhla, Maroc.

L'application offre aux clients un portail moderne pour d\'ecouvrir et r\'eserver des
circuits, s\'ejours et excursions, et expose aux gestionnaires un backoffice complet
incluant un programme de fid\'elit\'e, un calendrier de disponibilit\'e, la gestion
des ressources logistiques (guides, h\^otels, transports) et des statistiques en
temps r\'eel.

\bigskip
\textbf{Technologies cl\'es :}
\begin{itemize}
  \item \textbf{CMS :} WordPress 6.x avec WooCommerce 10.x
  \item \textbf{Backend :} PHP 8.x, plugin custom \texttt{travel-maroc-booking}
        (architecture MVC avec 8 classes sp\'ecialis\'ees)
  \item \textbf{Base de donn\'ees :} MySQL avec 10 tables m\'etier personnalis\'ees
  \item \textbf{Th\`eme :} Th\`eme WordPress sur mesure \texttt{travel-maroc}
        avec templates WooCommerce override
  \item \textbf{Multilingue :} Polylang 3.8 (Fran\c{c}ais / Arabe RTL)
  \item \textbf{Cartographie :} Leaflet.js + OpenStreetMap pour la carte des destinations
  \item \textbf{S\'ecurit\'e :} Nonces WP, capabilities, sanitisation, journal d'audit
  \item \textbf{SEO :} Open Graph, Twitter Cards, JSON-LD TouristTrip, Sitemap natif
\end{itemize}

\bigskip
\textbf{Mots-cl\'es :}
WordPress, WooCommerce, PHP, MySQL, Plugin custom, Programme fid\'elit\'e,
Multilingue, RTL, Leaflet.js, SEO.

\clearpage

\chapter*{Liste des abr\'eviations}
\addcontentsline{toc}{chapter}{Liste des abr\'eviations}
\thispagestyle{fancy}

\begin{table}[H]
\centering
\begin{tabularx}{0.9\linewidth}{>{\bfseries}lX}
  \toprule
  Abr\'eviation & Signification \\
  \midrule
  CMS     & Content Management System \\
  WC      & WooCommerce \\
  TMA     & Travel Maroc Agency \\
  TMB     & Travel Maroc Booking (pr\'efixe du plugin) \\
  CRUD    & Create, Read, Update, Delete \\
  ORM     & Object-Relational Mapping \\
  MCD     & Mod\`ele Conceptuel de Donn\'ees \\
  API     & Application Programming Interface \\
  SEO     & Search Engine Optimization \\
  OG      & Open Graph (protocole de m\'etadonn\'ees) \\
  RTL     & Right to Left (sens d'\'ecriture arabe) \\
  GA4     & Google Analytics 4 \\
  JWT     & JSON Web Token \\
  AJAX    & Asynchronous JavaScript and XML \\
  WP      & WordPress \\
  \bottomrule
\end{tabularx}
\end{table}

\clearpage
\tableofcontents
\clearpage
\listoftables
\addcontentsline{toc}{chapter}{Liste des tableaux}
\clearpage
\listoffigures
\addcontentsline{toc}{chapter}{Liste des figures}
\clearpage

% ─────────────────────────────────────────────────────────────
%  CORPS
% ─────────────────────────────────────────────────────────────
\cleardoublepage
\pagenumbering{arabic}

% ───────────────────────────────────────────────────────
\chapter{Contexte et Probl\'ematique}
% ───────────────────────────────────────────────────────

\section{Contexte G\'en\'eral}

Travel Maroc Agency est une agence de voyage fond\'ee \`a Dakhla en 2020.
Elle propose des circuits nationaux (d\'esert de Merzouga, Chefchaouen,
Marrakech), des s\'ejours internationaux (Istanbul, Duba\"i, Paris) et des
excursions \`a la journ\'ee. L'agence dispose de guides certifi\'es, de partenaires
h\^oteliers s\'electionn\'es et de compagnies de transport r\'eput\'ees.

Avant ce projet, la gestion des r\'eservations reposait enti\`erement sur des
\'echanges t\'el\'ephoniques et WhatsApp, sans outil num\'erique centralis\'e.

\begin{table}[H]
\centering
\caption{Probl\`emes identifi\'es dans la gestion actuelle}
\begin{tabularx}{\textwidth}{>{\bfseries}lX}
  \toprule
  Probl\`eme & Impact \\
  \midrule
  Catalogue de voyages non structur\'e &
    Clients ne peuvent pas comparer les offres facilement \\
  R\'eservations uniquement par t\'el\'ephone &
    Disponibilit\'e limit\'ee, erreurs de saisie, pas de trace \\
  Aucun programme de fid\'elit\'e &
    Perte de clients r\'ecurrents face \`a la concurrence \\
  Gestion manuelle des ressources &
    Guides, h\^otels et transports non rattach\'es aux offres \\
  Pas de pr\'esence web multilingue &
    Clientele arabe et internationale non adress\'ee \\
  \bottomrule
\end{tabularx}
\end{table}

\section{Probl\'ematique}

\begin{infobox}
\textbf{Comment} concevoir et d\'evelopper une plateforme web de r\'eservation
de voyages permettant \`a Travel Maroc Agency de digitaliser son catalogue,
d'automatiser la gestion des r\'eservations, de fid\'eliser ses clients via un
programme de points, et de donner aux gestionnaires un tableau de bord
op\'erationnel complet~?
\end{infobox}

\section{P\'erim\`etre du Projet}

Ce projet s'inscrit dans le cadre du module \textbf{D\'eveloppement Web Avanc\'e}
et couvre~:

\begin{itemize}
  \item D\'eveloppement d'un CMS WordPress avec WooCommerce comme moteur e-commerce
  \item Cr\'eation d'un plugin m\'etier PHP orienté objet (8 classes)
  \item Conception d'un th\`eme WordPress sur mesure avec templates WooCommerce
  \item Mod\'elisation et cr\'eation de 10 tables MySQL sp\'ecialis\'ees
  \item Int\'egration du multilingue (Fran\c{c}ais / Arabe) avec support RTL
  \item S\'ecurisation et journalisation de toutes les actions d'administration
\end{itemize}

% ───────────────────────────────────────────────────────
\chapter{Objectifs du Projet}
% ───────────────────────────────────────────────────────

\section{Objectifs Fonctionnels}

\begin{enumerate}
  \item \textbf{Catalogue en ligne} : Affichage des offres de voyage avec
        destination, dur\'ee, points fid\'elit\'e, guide, h\^otel et transport associ\'es
  \item \textbf{R\'eservation en ligne} : Processus de checkout WooCommerce \'etendu
        avec saisie de la date de d\'epart, nombre de voyageurs et demandes sp\'eciales
  \item \textbf{Calendrier de disponibilit\'e} : Blocage de dates par offre depuis
        le backoffice, avec validation c\^ot\'e serveur au checkout
  \item \textbf{Programme de fid\'elit\'e} : Attribution de points apr\`es chaque
        r\'eservation confirm\'ee, remise automatique selon le palier client
        (Standard / Silver / Gold / VIP)
  \item \textbf{Espace client enrichi} : Tableau de bord avec solde de points,
        notifications et historique de r\'eservations avec d\'etails voyage
  \item \textbf{Backoffice m\'etier} : CRUD Destinations, Guides, H\^otels, Transport
  \item \textbf{Dashboard administrateur} : KPIs en temps r\'eel et export CSV des
        r\'eservations
  \item \textbf{Multilingue Fran\c{c}ais / Arabe} : Sélecteur de langue dans le
        header avec support complet RTL pour l'arabe
\end{enumerate}

\section{Objectifs Techniques}

\begin{enumerate}
  \item Ma\^itriser l'architecture de plugin WordPress (hooks, filters, classes)
  \item Impl\'ementer un syst\`eme de base de donn\'ees relationnelle sur mesure (MySQL)
  \item Appliquer les bonnes pratiques de s\'ecurit\'e WordPress (nonces, capabilities,
        sanitisation, \'echappement)
  \item Int\'egrer des librairies JavaScript modernes (Leaflet.js, GSAP)
  \item Produire des m\'etadonn\'ees SEO automatiques (Open Graph, JSON-LD)
\end{enumerate}

% ───────────────────────────────────────────────────────
\chapter{Environnement de D\'eveloppement}
% ───────────────────────────────────────────────────────

\section{Logiciels et Outils}

\begin{table}[H]
\centering
\caption{Logiciels et outils utilis\'es}
\begin{tabularx}{\linewidth}{llX}
  \toprule
  \textbf{Outil} & \textbf{Version} & \textbf{R\^ole} \\
  \midrule
  XAMPP          & v8.2.x     & Serveur local (Apache + MySQL + PHP) \\
  PHP            & v8.2.x     & Langage serveur principal \\
  WordPress      & v6.7.x     & CMS de base \\
  WooCommerce    & v10.7.x    & Moteur e-commerce \\
  Polylang       & v3.8.x     & Plugin multilingue \\
  MySQL          & v8.0.x     & Base de donn\'ees relationnelle \\
  phpMyAdmin     & v5.x       & Administration visuelle de la BDD \\
  Visual Studio Code & v1.9x  & \'Editeur de code principal \\
  Git            & v2.x       & Contr\^ole de version \\
  Postman        & v11.x      & Test des endpoints REST TMA \\
  \bottomrule
\end{tabularx}
\end{table}

\section{Stack Technique}

\begin{table}[H]
\centering
\caption{Technologies utilis\'ees par couche}
\begin{tabularx}{\linewidth}{lX}
  \toprule
  \textbf{Couche} & \textbf{Technologies} \\
  \midrule
  Serveur         & Apache 2.4, PHP 8.2 (XAMPP) \\
  CMS             & WordPress 6.7 + WooCommerce 10.7 \\
  Plugin m\'etier & PHP 8 orienté objet, 8 classes (pattern Singleton/Static) \\
  Th\`eme         & PHP, HTML5, CSS3 personnalis\'e, WordPress Template Hierarchy \\
  Base de donn\'ees & MySQL 8 (10 tables custom via \texttt{dbDelta}) \\
  Frontend        & JavaScript ES6+, GSAP 3.12, Leaflet.js 1.9, jQuery (WC) \\
  Multilingue     & Polylang 3.8 (API native : \texttt{pll\_\_()}, \texttt{pll\_register\_string()}) \\
  SEO             & Open Graph, Twitter Cards, Schema.org JSON-LD, Sitemap natif \\
  Analytics       & Google Analytics 4 (snippet GA4 inject\'e via \texttt{wp\_head}) \\
  \bottomrule
\end{tabularx}
\end{table}

\section{Structure du Projet}

\begin{shellbox}
c:\xampp\htdocs\boutique\
+-- wp-content/
    +-- plugins/
    |   +-- travel-maroc-booking/           # Plugin metier principal
    |       +-- travel-maroc-booking.php    # Point d'entree
    |       +-- includes/
    |       |   +-- class-tmb-activator.php    # Activation, tables, seed
    |       |   +-- class-tmb-fidelite.php     # Programme de points
    |       |   +-- class-tmb-checkout.php     # Champs voyage au checkout
    |       |   +-- class-tmb-availability.php # Calendrier disponibilite
    |       |   +-- class-tmb-notifications.php
    |       |   +-- class-tmb-roles.php        # Roles et permissions
    |       |   +-- class-tmb-rest-api.php     # Endpoints REST
    |       |   +-- class-tmb-log-admin.php    # Journal d'audit
    |       +-- admin/
    |           +-- class-tmb-admin.php        # Dashboard + export CSV
    |           +-- class-tmb-resources-admin.php # CRUD Guides/Hotels/Transport
    |           +-- class-tmb-product-meta.php # Meta box produit
    |           +-- class-tmb-settings.php     # Page reglages
    +-- themes/
        +-- travel-maroc/                   # Theme custom
            +-- functions.php               # Hooks, SEO, Polylang, GA4
            +-- header.php  footer.php
            +-- assets/css/main.css
            +-- assets/js/main.js
            +-- template-parts/             # about, contact, offer-card
            +-- woocommerce/                # Overrides templates WC
            |   +-- single-product.php
            |   +-- myaccount/dashboard.php orders.php
            |   +-- emails/
            +-- wpml-config.xml             # Sync Polylang/WPML
\end{shellbox}

% ───────────────────────────────────────────────────────
\chapter{Architecture G\'en\'erale du Syst\`eme}
% ───────────────────────────────────────────────────────

\section{Vue d'Ensemble --- Architecture WordPress}

L'application repose sur l'architecture \textbf{WordPress Hook System}~: le plugin
et le th\`eme communiquent exclusivement via des \textit{actions} et des
\textit{filtres}, sans couplage direct entre les classes.

\begin{figure}[H]
\centering
\diagramtitle{Architecture applicative en trois couches}
\begin{tikzpicture}[node distance=1.45cm]
  \node[diagram box, text width=11cm] (front) {
    {\bfseries\color{estblue}Couche Pr\'esentation (Th\`eme)}\\[3pt]
    Templates PHP, CSS sur mesure, JavaScript (GSAP, Leaflet),
    overrides WooCommerce, Polylang RTL
  };
  \node[diagram box gold, below=of front, text width=11cm] (plug) {
    {\bfseries\color{estblue}Couche M\'etier (Plugin TMB)}\\[3pt]
    Fid\'elit\'e, Checkout, Disponibilit\'es, Notifications, Rôles,
    API REST, Journal d'audit, Param\`etres
  };
  \node[diagram db, below=of plug, text width=11cm] (db) {
    {\bfseries\color{estblue}Couche Donn\'ees}\\[3pt]
    MySQL~: tables WordPress standard + 10 tables m\'etier TMA
    (\texttt{tma\_client\_profile}, \texttt{tma\_historique\_points}, \ldots)
  };
  \draw[diagram arrow] (front) -- node[diagram note, right] {Hooks WP (add\_action / apply\_filters)} (plug);
  \draw[diagram arrow] (plug)  -- node[diagram note, right] {\texttt{\$wpdb} + dbDelta} (db);
  \begin{scope}[on background layer]
    \node[diagram frame, fit=(front)(plug)(db), inner sep=12pt, fill=estblue!2] {};
  \end{scope}
\end{tikzpicture}
\caption{Architecture client-serveur \`a trois couches}
\end{figure}

\section{Acteurs et Rôles}

\begin{table}[H]
\centering
\caption{Rôles et permissions dans le syst\`eme}
\begin{tabularx}{\linewidth}{lXl}
  \toprule
  \textbf{Rôle} & \textbf{Acc\`es} & \textbf{Capacit\'e WordPress} \\
  \midrule
  Visiteur      & Pages publiques, catalogue, formulaire contact & Aucune \\
  Client        & Mon compte, checkout, fid\'elit\'e, notifications & \texttt{read} \\
  Agent TMA     & Dashboard TMA (lecture), clients, commandes    & \texttt{tma\_view\_*} \\
  Shop Manager  & WC complet + CRUD m\'etier + export CSV         & \texttt{manage\_woocommerce} \\
  Administrateur & Tout + r\'eglages + journal d'audit            & \texttt{manage\_options} \\
  \bottomrule
\end{tabularx}
\end{table}

\section{Flux d'une R\'eservation}

\begin{figure}[H]
\centering
\diagramtitle{Sc\'enario complet de r\'eservation d'un voyage}
\begin{tikzpicture}[
  scale=0.8, transform shape,
  lifeline/.style={draw=gray!65, line width=0.9pt, dashed},
  msg/.style={-{Stealth[length=2.6mm,width=2mm]}, estblue, line width=1.1pt},
  returnmsg/.style={-{Stealth[length=2.6mm,width=2mm]}, estgold!70!black, dashed, line width=1pt},
  label/.style={font=\scriptsize, fill=white, inner sep=1.5pt, text=gray!85}
]
  \node[diagram actor]   (client) at (0,0)    {Client};
  \node[diagram box, minimum width=3.2cm, text width=3.2cm] (theme)  at (4.2,0)   {Th\`eme WC};
  \node[diagram box gold, minimum width=3.2cm, text width=3.2cm] (plugin) at (8.6,0)   {Plugin TMB};
  \node[diagram db, minimum width=3cm, text width=3cm]       (db)     at (13.0,0)  {MySQL};

  \draw[lifeline] (client.south) -- ++(0,-9cm);
  \draw[lifeline] (theme.south)  -- ++(0,-9cm);
  \draw[lifeline] (plugin.south) -- ++(0,-9cm);
  \draw[lifeline] (db.south)     -- ++(0,-9cm);

  \draw[msg]       (client.south)++(0,-0.5) -- node[label,above]{S\'electionne une offre}       ++(4.2,0);
  \draw[msg]       (theme.south) ++(0,-1.3) -- node[label,above]{Checkout : date, voyageurs}    ++(4.4,0);
  \draw[msg]       (plugin.south)++(0,-2.1) -- node[label,above]{V\'erifie disponibilit\'e}     ++(4.4,0);
  \draw[returnmsg] (db.south)    ++(0,-2.9) -- node[label,above]{Date non bloqu\'ee}             ++(-4.4,0);
  \draw[msg]       (plugin.south)++(0,-3.7) -- node[label,above]{V\'erifie remise fid\'elit\'e} ++(4.4,0);
  \draw[returnmsg] (db.south)    ++(0,-4.5) -- node[label,above]{Taux palier client}             ++(-4.4,0);
  \draw[returnmsg] (plugin.south)++(0,-5.3) -- node[label,above]{Applique remise au panier}     ++(-4.4,0);
  \draw[msg]       (client.south)++(0,-6.1) -- node[label,above]{Confirme la commande}           ++(4.2,0);
  \draw[msg]       (plugin.south)++(0,-6.9) -- node[label,above]{Sauvegarde m\'etadonn\'ees}    ++(4.4,0);
  \draw[returnmsg] (db.south)    ++(0,-7.7) -- node[label,above]{Commande persist\'ee}           ++(-4.4,0);
  \draw[returnmsg] (client.south)++(0,-8.5) -- node[label,above]{Email + notification}           ++(-4.2,0);
\end{tikzpicture}
\caption{Diagramme de s\'equence : r\'eservation d'une offre de voyage}
\end{figure}

% ───────────────────────────────────────────────────────
\chapter{Architecture D\'etaill\'ee (C4 et Dynamique)}
% ───────────────────────────────────────────────────────

\section{Diagramme C4 --- Niveau 1 : Contexte Syst\`eme}

\begin{figure}[H]
\centering
\diagramtitle{C4 -- Niveau 1 : vision contexte}
\begin{tikzpicture}[x=1cm,y=1cm,scale=0.90,transform shape]
  \node[diagram actor] (client)  at (0,1.5)  {Client};
  \node[diagram actor] (manager) at (0,-1.5) {Gestionnaire};
  \node[diagram box, minimum width=6cm, text width=6cm] (tma) at (6,0) {
    \textbf{Travel Maroc Agency}\\Catalogue, r\'eservation, fid\'elit\'e, backoffice
  };
  \node[diagram db, minimum width=3.5cm, text width=3.5cm] (mysql) at (11.5,-1.8) {
    MySQL\\10 tables m\'etier
  };
  \node[diagram frame, fill=gray!8, minimum width=3.5cm, minimum height=1.3cm,
        align=center, font=\small] (wp) at (11.5,1.8) {
    WordPress\\WooCommerce
  };
  \draw[diagram arrow] (client)  -- node[diagram note,above]{R\'eserve des voyages}     (tma);
  \draw[diagram arrow] (manager) -- node[diagram note,below]{Administre le contenu}     (tma);
  \draw[diagram arrow] (tma)     -- node[diagram note,above]{Hooks + \$wpdb}            (wp);
  \draw[diagram arrow] (tma)     -- node[diagram note,right]{dbDelta + CRUD}            (mysql);
  \begin{scope}[on background layer]
    \node[diagram frame, fit=(client)(manager)(tma)(wp)(mysql), inner sep=14pt, fill=estblue!2] {};
  \end{scope}
\end{tikzpicture}
\caption{C4 Niveau 1 : Contexte du syst\`eme Travel Maroc Agency}
\end{figure}

\section{Diagramme C4 --- Niveau 2 : Conteneurs}

\begin{figure}[H]
\centering
\diagramtitle{C4 -- Niveau 2 : conteneurs applicatifs}
\begin{tikzpicture}[x=1cm,y=1cm,scale=0.82,transform shape]
  \node[diagram actor] (cli) at (0,1.6)  {Client};
  \node[diagram actor] (adm) at (0,-1.6) {Admin};
  \node[diagram box, minimum width=4.8cm, text width=4.8cm] (theme) at (5,1.6) {
    \textbf{Th\`eme travel-maroc}\\Templates, CSS, JS, Overrides WC
  };
  \node[diagram box gold, minimum width=4.8cm, text width=4.8cm] (plugin) at (10.5,1.6) {
    \textbf{Plugin TMB}\\8 classes PHP, hooks WP, REST API
  };
  \node[diagram db, minimum width=4.8cm, text width=4.8cm] (db) at (10.5,-1.6) {
    \textbf{MySQL}\\wp\_* (WP/WC) + tma\_* (m\'etier)
  };
  \draw[diagram arrow] (cli.east) -- (theme.west);
  \draw[diagram arrow] (adm.east) -- ++(1,0) |- (plugin.west);
  \draw[diagram arrow] (theme)    -- node[diagram note,above]{add\_action / filters} (plugin);
  \draw[diagram arrow] (plugin)   -- node[diagram note,right]{\$wpdb + dbDelta}      (db);
\end{tikzpicture}
\caption{C4 Niveau 2 : Conteneurs (th\`eme, plugin, BDD)}
\end{figure}

\section{Diagramme C4 --- Niveau 3 : Composants du Plugin}

\begin{figure}[H]
\centering
\diagramtitle{D\'ecomposition interne du plugin travel-maroc-booking}
\begin{tikzpicture}[x=1cm,y=1cm,scale=0.76,transform shape]
  \node[diagram box, minimum width=3.5cm, text width=3.5cm] (act) at (0,3.5) {
    \textbf{TMB\_Activator}\\Tables, seed, crons};
  \node[diagram box, minimum width=3.5cm, text width=3.5cm] (fid) at (4.5,3.5) {
    \textbf{TMB\_Fidelite}\\Points, paliers, remises};
  \node[diagram box, minimum width=3.5cm, text width=3.5cm] (co)  at (9,3.5) {
    \textbf{TMB\_Checkout}\\Champs voyage, validation};
  \node[diagram box, minimum width=3.5cm, text width=3.5cm] (av)  at (13.5,3.5) {
    \textbf{TMB\_Availability}\\Calendrier, blocage dates};
  \node[diagram box gold, minimum width=3.5cm, text width=3.5cm] (adm) at (0,1) {
    \textbf{TMB\_Admin}\\Dashboard, export CSV};
  \node[diagram box gold, minimum width=3.5cm, text width=3.5cm] (res) at (4.5,1) {
    \textbf{TMB\_Resources}\\CRUD G/H/T};
  \node[diagram box gold, minimum width=3.5cm, text width=3.5cm] (set) at (9,1) {
    \textbf{TMB\_Settings}\\R\'eglages, comptes};
  \node[diagram box, minimum width=3.5cm, text width=3.5cm] (log) at (13.5,1) {
    \textbf{TMB\_Log\_Admin}\\Journal d'audit};
  \node[diagram db, minimum width=14cm, text width=14cm] (db) at (6.75,-1.2) {
    \textbf{MySQL} --- \$wpdb --- 10 tables m\'etier TMA
  };
  \foreach \s in {act,fid,co,av,adm,res,set,log} {
    \draw[diagram arrow] (\s.south) -- ++(0,-0.3) |- (db.north);
  }
\end{tikzpicture}
\caption{C4 Niveau 3 : Composants internes du plugin}
\end{figure}

\section{Architecture du Th\`eme}

\begin{table}[H]
\centering
\caption{Couches du th\`eme travel-maroc}
\begin{tabularx}{\textwidth}{lX}
  \toprule
  \textbf{Fichier / Dossier} & \textbf{Responsabilit\'e} \\
  \midrule
  \texttt{functions.php}   & Enregistrement des hooks, SEO OG, GA4, Polylang,
                             sélecteur de langue, RTL, sitemap \\
  \texttt{header.php}      & Navigation, logo GSAP animé, panier, sélecteur langue \\
  \texttt{footer.php}      & Liens sociaux, copyright \\
  \texttt{assets/css/main.css} & Styles sur mesure + variables CSS + RTL arabe \\
  \texttt{assets/js/main.js}   & Interactions, formulaire AJAX contact \\
  \texttt{assets/js/logo-animation.js} & Animation SVG logo via GSAP 3 \\
  \texttt{woocommerce/single-product.php} & Fiche produit avec ressources TMA \\
  \texttt{woocommerce/myaccount/orders.php} & Réservations enrichies \\
  \texttt{woocommerce/emails/}  & Emails personnalis\'es (processing, completed) \\
  \texttt{template-parts/}      & Blocs r\'eutilisables (carte offre, about, contact) \\
  \bottomrule
\end{tabularx}
\end{table}

% ───────────────────────────────────────────────────────
\chapter{Mod\'elisation des Donn\'ees}
% ───────────────────────────────────────────────────────

\section{Diagramme MCD --- Tables M\'etier TMA}

En plus des tables standard WordPress (\texttt{wp\_posts}, \texttt{wp\_users},
\texttt{wp\_postmeta}\ldots) et WooCommerce, le plugin cr\'ee 10 tables
personnalis\'ees via \texttt{dbDelta()}~:

\begin{figure}[H]
\centering
\diagramtitle{Relations entre les 10 tables m\'etier TMA}
\begin{tikzpicture}[x=1cm,y=1cm,scale=0.78,transform shape,
  tbl/.style={diagram frame, fill=lightblue!30, text width=3.6cm,
    align=left, inner sep=8pt, font=\tiny}]

  \node[tbl] (loc) at (0,6) {
    \textbf{tma\_localisation}\\
    \textcolor{estgold}{\#id}\\nom, type, code\_pays\\latitude, longitude, actif
  };
  \node[tbl] (dest) at (5,6) {
    \textbf{tma\_destination}\\
    \textcolor{estgold}{\#id}, localisation\_id\\nom, pays, region\\description, statut
  };
  \node[tbl] (guide) at (0,3) {
    \textbf{tma\_guide}\\
    \textcolor{estgold}{\#id}, localisation\_id\\prenom, nom, telephone\\langues, tarif\_journalier
  };
  \node[tbl] (hotel) at (5,3) {
    \textbf{tma\_hotel}\\
    \textcolor{estgold}{\#id}, localisation\_id\\nom, categorie\_etoiles\\prix\_chambre\_nuit
  };
  \node[tbl] (trans) at (10,3) {
    \textbf{tma\_transport}\\
    \textcolor{estgold}{\#id}\\localisation\_depart\_id\\localisation\_arrivee\_id\\type, compagnie
  };
  \node[tbl] (tc) at (10,6) {
    \textbf{tma\_type\_client}\\
    \textcolor{estgold}{\#id}\\libelle, taux\_remise\\seuil\_points\\points\_jamais\_expirent
  };
  \node[tbl] (prof) at (7,0.5) {
    \textbf{tma\_client\_profile}\\
    \textcolor{estgold}{\#id}, wp\_user\_id\\type\_client\_id\\points\_fidelite\_solde
  };
  \node[tbl] (hist) at (3,0.5) {
    \textbf{tma\_historique\_points}\\
    \textcolor{estgold}{\#id}, client\_id\\wc\_order\_id\\operation, points\\date\_expiration
  };
  \node[tbl] (notif) at (10,0) {
    \textbf{tma\_notification}\\
    \textcolor{estgold}{\#id}, client\_id\\type, titre, message\\statut, canal
  };
  \node[tbl] (log) at (0,0.5) {
    \textbf{tma\_log\_admin}\\
    \textcolor{estgold}{\#id}, wp\_user\_id\\action, entite\_cible\\ip\_address
  };

  \draw[diagram arrow] (dest.west)  -- (loc.east);
  \draw[diagram arrow] (guide.north)-- (loc.south);
  \draw[diagram arrow] (hotel.north)-- (loc.south);
  \draw[diagram arrow] (trans.north)-- (loc.south);
  \draw[diagram arrow] (prof.north) -- (tc.south);
  \draw[diagram arrow] (hist.east)  -- (prof.west);
  \draw[diagram arrow] (notif.west) -- (prof.east);
\end{tikzpicture}
\caption{MCD des tables m\'etier TMA}
\end{figure}

\section{Tables Principales}

\begin{table}[H]
\centering
\caption{Table \texttt{tma\_type\_client} --- Paliers de fid\'elit\'e}
\begin{tabularx}{\linewidth}{llX}
  \toprule
  \textbf{Colonne} & \textbf{Type} & \textbf{Description} \\
  \midrule
  \texttt{id}                     & INT        & Cl\'e primaire \\
  \texttt{libelle}                & VARCHAR(50) & Standard / Silver / Gold / VIP / Entreprise \\
  \texttt{taux\_remise}           & DECIMAL(5,2) & Pourcentage de remise automatique (0--20) \\
  \texttt{seuil\_points}          & INT         & Points cumul\'es pour atteindre ce palier \\
  \texttt{points\_jamais\_expirent} & TINYINT   & 1 = points permanents (VIP uniquement) \\
  \texttt{couleur\_badge}         & VARCHAR(7)  & Code hexad\'ecimal pour l'affichage \\
  \bottomrule
\end{tabularx}
\end{table}

\begin{table}[H]
\centering
\caption{Table \texttt{tma\_historique\_points} --- Transactions fid\'elit\'e}
\begin{tabularx}{\linewidth}{llX}
  \toprule
  \textbf{Colonne} & \textbf{Type} & \textbf{Description} \\
  \midrule
  \texttt{id}             & INT           & Cl\'e primaire \\
  \texttt{client\_id}     & INT           & R\'ef. \texttt{tma\_client\_profile.id} \\
  \texttt{wc\_order\_id}  & BIGINT        & R\'ef. commande WooCommerce \\
  \texttt{operation}      & VARCHAR(20)   & GAIN / ANNULATION / EXPIRATION / DEDUCTION \\
  \texttt{points}         & INT           & Valeur (positive ou n\'egative) \\
  \texttt{solde\_avant}   & INT UNSIGNED  & Solde avant op\'eration \\
  \texttt{solde\_apres}   & INT UNSIGNED  & Solde apr\`es op\'eration \\
  \texttt{date\_expiration} & DATE        & Date d'expiration du GAIN (null si VIP) \\
  \bottomrule
\end{tabularx}
\end{table}

\begin{table}[H]
\centering
\caption{Table \texttt{tma\_notification} --- File de notifications}
\begin{tabularx}{\linewidth}{llX}
  \toprule
  \textbf{Colonne} & \textbf{Type} & \textbf{Description} \\
  \midrule
  \texttt{type}   & VARCHAR(40) & CONFIRMATION\_RESERVATION / POINTS\_GAGNES / MONTEE\_PALIER / ANNULATION \\
  \texttt{canal}  & VARCHAR(10) & EMAIL / IN\_APP \\
  \texttt{statut} & VARCHAR(15) & EN\_ATTENTE / ENVOYE / ERREUR / LU \\
  \bottomrule
\end{tabularx}
\end{table}

% ───────────────────────────────────────────────────────
\chapter{Impl\'ementation --- Plugin M\'etier}
% ───────────────────────────────────────────────────────

\section{Programme de Fid\'elit\'e --- \texttt{TMB\_Fidelite}}

La classe \texttt{TMB\_Fidelite} g\`ere l'int\'egralit\'e du cycle de vie des points~:
attribution, annulation, expiration et remise automatique au panier.

\subsection{R\`egles de calcul}

\begin{table}[H]
\centering
\caption{Param\`etres du programme de fid\'elit\'e}
\begin{tabularx}{\linewidth}{lX}
  \toprule
  \textbf{R\`egle} & \textbf{Valeur} \\
  \midrule
  Taux de gain     & 0,1 point par MAD d\'epens\'e (1 point = 10 MAD) \\
  Valeur du point  & 100 points = 10 MAD de remise \\
  Expiration       & 24 mois apr\`es la date de gain \\
  Exception VIP    & Points sans expiration (\texttt{points\_jamais\_expirent = 1}) \\
  Palier Silver    & 500 points cumul\'es $\Rightarrow$ remise 5\% \\
  Palier Gold      & 1~500 points cumul\'es $\Rightarrow$ remise 10\% \\
  Palier VIP       & 5~000 points cumul\'es $\Rightarrow$ remise 20\% \\
  \bottomrule
\end{tabularx}
\end{table}

\subsection{Attribution automatique des points}

\begin{phpbox}
// Hook sur woocommerce_order_status_completed
public static function attribuer_points( int $order_id ): void {
    $order   = wc_get_order( $order_id );
    $profile = self::get_profile( $order->get_user_id() );
    $total   = (float) $order->get_total();
    $points  = (int) floor( $total * self::POINTS_PAR_MAD ); // 0.1 pt/MAD

    $exp = $profile->points_jamais_expirent
        ? null
        : gmdate( 'Y-m-d', strtotime( '+24 months' ) );

    self::crediter( $profile->id, $points, $order_id,
        "Commande #$order_id", $exp );
    self::verifier_montee_palier( $profile->id );
}
\end{phpbox}

\subsection{Remise automatique au panier}

\begin{phpbox}
// Hook sur woocommerce_cart_calculate_fees
public static function appliquer_remise_panier( WC_Cart $cart ): void {
    $user_id = get_current_user_id();
    if ( ! $user_id ) return;

    $profile = self::get_profile( $user_id );
    if ( ! $profile || ! $profile->taux_remise ) return;

    $total   = $cart->get_subtotal();
    $remise  = round( $total * ( $profile->taux_remise / 100 ), 2 );
    $cart->add_fee( self::FEE_LABEL, -$remise ); // valeur negative
}
\end{phpbox}

\section{Champs Voyage au Checkout --- \texttt{TMB\_Checkout}}

\begin{phpbox}
// Hook sur woocommerce_after_order_notes
public static function add_fields( WC_Checkout $checkout ): void {
    woocommerce_form_field( 'tma_date_depart', [
        'type'     => 'text',
        'label'    => 'Date de depart souhaitee',
        'required' => true,
    ], $checkout->get_value('tma_date_depart') );

    woocommerce_form_field( 'tma_nb_adultes', [
        'type'     => 'number',
        'label'    => 'Adultes (12 ans et +)',
        'required' => true,
        'default'  => 1,
    ], $checkout->get_value('tma_nb_adultes') ?: 1 );
}

// Hook sur woocommerce_checkout_process (validation serveur)
public static function validate_fields(): void {
    $date = sanitize_text_field( $_POST['tma_date_depart'] ?? '' );
    if ( ! $date ) {
        wc_add_notice( 'Date de depart obligatoire.', 'error' );
    }
    // Verification dates bloquees pour chaque produit du panier
    foreach ( WC()->cart->get_cart() as $item ) {
        if ( ! TMB_Availability::is_date_available( $item['product_id'], $date ) ) {
            wc_add_notice( 'Cette date est indisponible.', 'error' );
        }
    }
}
\end{phpbox}

\section{Calendrier de Disponibilit\'e --- \texttt{TMB\_Availability}}

Le calendrier de disponibilit\'e est impl\'ement\'e comme une \textit{meta box}
native WordPress sur chaque fiche produit (offre de voyage).

\begin{itemize}
  \item \textbf{Stockage} : Dates bloqu\'ees en JSON dans \texttt{\_tma\_dates\_bloquees}
  \item \textbf{Interface} : Calendrier JavaScript cliquable (mois/ann\'ee, fond rouge = bloqu\'e)
  \item \textbf{Validation} : \texttt{TMB\_Availability::is\_date\_available()} appel\'e
        lors de la validation du checkout
  \item \textbf{Capacit\'e} : Option \texttt{\_tma\_capacite\_max} pour d\'efinir le
        nombre maximal de voyageurs par d\'epart (0 = illimit\'e)
\end{itemize}

\section{Journal d'Audit --- \texttt{TMB\_Log\_Admin}}

Toutes les actions du backoffice sont trac\'ees dans \texttt{tma\_log\_admin}~:

\begin{phpbox}
public static function write( array $data ): void {
    global $wpdb;
    $wpdb->insert( $wpdb->prefix . 'tma_log_admin', [
        'wp_user_id'   => get_current_user_id(),
        'action'       => $data['action'],        // CREATE/UPDATE/DELETE
        'entite_cible' => $data['entite_cible'],  // tma_destination, tma_guide...
        'entite_id'    => $data['entite_id'] ?? null,
        'ip_address'   => $_SERVER['REMOTE_ADDR'], // REMOTE_ADDR uniquement (anti-spoofing)
        'user_agent'   => substr( $_SERVER['HTTP_USER_AGENT'] ?? '', 0, 255 ),
        'date_action'  => current_time( 'mysql' ),
    ]);
}
\end{phpbox}

\begin{notebox}
\textbf{S\'ecurit\'e IP} : L'adresse IP est lue exclusivement depuis
\texttt{\$\_SERVER['REMOTE\_ADDR']}. Les en-t\^etes \texttt{X-Forwarded-For}
et \texttt{X-Real-IP} sont ignor\'es pour emp\^echer le \textit{IP spoofing}.
\end{notebox}

% ───────────────────────────────────────────────────────
\chapter{Impl\'ementation --- Th\`eme et Frontend}
% ───────────────────────────────────────────────────────

\section{SEO --- Open Graph et Donn\'ees Structur\'ees}

\subsection{Balises Open Graph}

\begin{phpbox}
function tma_og_meta(): void {
    global $post;
    if ( is_singular('product') && $post ) {
        $product = wc_get_product( $post->ID );
        $title   = $product->get_name();
        $desc    = wp_trim_words( $product->get_description(), 25 );
        $img     = get_the_post_thumbnail_url( $post->ID, 'large' );
        $dest    = get_post_meta( $post->ID, '_tma_destination', true );
        if ( $dest ) $desc = $dest . ' -- ' . $desc;
    }
    echo '<meta property="og:title" content="' . esc_attr($title) . '">';
    echo '<meta property="og:description" content="' . esc_attr($desc) . '">';
    echo '<meta name="twitter:card" content="summary_large_image">';
}
add_action( 'wp_head', 'tma_og_meta' );
\end{phpbox}

\subsection{Donn\'ees Structur\'ees JSON-LD}

Chaque fiche produit g\'en\`ere automatiquement un bloc \texttt{TouristTrip}
reconnu par Google~:

\begin{jsonbox}
{
  "@context": "https://schema.org",
  "@type": "TouristTrip",
  "name": "Circuit Desert Merzouga 7 jours",
  "description": "Experience unique dans les dunes de sable dore...",
  "offers": {
    "@type": "Offer",
    "price": 2500,
    "priceCurrency": "MAD",
    "availability": "https://schema.org/InStock"
  }
}
\end{jsonbox}

\section{Carte Interactive des Destinations --- Leaflet.js}

La page d'accueil affiche une carte interactive via le shortcode
\texttt{[tma\_destinations\_map]}, rendu c\^ot\'e client par Leaflet.js~:

\begin{itemize}
  \item Tuiles OpenStreetMap (sans cl\'e API, open-source)
  \item Marqueurs oranges sur chaque destination avec coordonn\'ees GPS
  \item Popup au clic~: nom, pays, description de la destination
  \item Seules les destinations avec latitude/longitude non nulles sont affich\'ees
\end{itemize}

\section{Animation Logo GSAP}

Le logo de l'agence est anim\'e via \textbf{GSAP 3.12} au chargement de la page~:

\begin{itemize}
  \item Phase 1~: Trac\'e progressif du path SVG (effet \textit{draw})
  \item Phase 2~: Fondu et rotation des \'el\'ements graphiques
  \item Phase 3~: Apparition du texte avec effet \textit{fade-in}
\end{itemize}

\section{Support RTL Arabe}

\begin{phpbox}
// Ajout de la classe tma-rtl quand la langue est RTL
add_filter( 'body_class', function( $classes ) {
    if ( get_bloginfo('text_direction') === 'rtl' ) {
        $classes[] = 'tma-rtl';
    }
    return $classes;
});
\end{phpbox}

La classe \texttt{.tma-rtl} active dans \texttt{main.css}~:
\begin{itemize}
  \item \texttt{direction: rtl} sur le body
  \item Police arabe : \texttt{Noto Naskh Arabic, serif}
  \item Inversion des flex-directions sur la navigation et les grilles
  \item Ajustement des marges et paddings
\end{itemize}

% ───────────────────────────────────────────────────────
\chapter{Fonctionnalit\'es Avanc\'ees}
% ───────────────────────────────────────────────────────

\section{Dashboard Administrateur avec KPIs Temps R\'eel}

Le tableau de bord TMA affiche des indicateurs calcul\'es directement
depuis la base de donn\'ees~:

\begin{table}[H]
\centering
\caption{KPIs affich\'es dans le dashboard TMA}
\begin{tabularx}{\linewidth}{lX}
  \toprule
  \textbf{KPI} & \textbf{Requ\^ete SQL} \\
  \midrule
  Revenus totaux      & \texttt{SUM(meta\_value)} sur \texttt{wp\_postmeta} où \texttt{meta\_key='\_order\_total'} et statut \texttt{wc-completed} \\
  Revenus ce mois     & Idem filtré sur \texttt{DATE\_FORMAT(post\_date, '\%Y-\%m')} \\
  Clients inscrits    & \texttt{COUNT(*)} dans \texttt{tma\_client\_profile} \\
  Points distribu\'es & \texttt{SUM(points)} dans \texttt{tma\_historique\_points} où \texttt{operation='GAIN'} \\
  Top 5 offres        & \texttt{GROUP BY product\_id ORDER BY COUNT(*) DESC} \\
  \bottomrule
\end{tabularx}
\end{table}

\section{Export CSV des R\'eservations}

Un bouton \textbf{``$\downarrow$ Exporter CSV r\'eservations''} est disponible dans le
dashboard et g\'en\`ere un fichier \texttt{reservations-tma-YYYY-MM-DD.csv}~:

\begin{phpbox}
header('Content-Type: text/csv; charset=utf-8');
header('Content-Disposition: attachment; filename="' . $filename . '"');
fwrite($out, "\xEF\xBB\xBF"); // BOM UTF-8 pour Excel

fputcsv($out, ['N° Commande','Date','Statut','Client','Email',
               'Offre','Total TTC','Date depart','Adultes',
               'Enfants','Demandes speciales'], ';');
\end{phpbox}

\begin{notebox}
Le fichier est encod\'e \textbf{UTF-8 avec BOM} pour assurer la compatibilit\'e avec
Microsoft Excel sans probl\`emes d'accents (en particulier pour les caract\`eres
arabes et fran\c{c}ais).
\end{notebox}

\section{API REST TMA}

Le plugin expose une API REST consommable par des applications tierces~:

\begin{table}[H]
\centering
\caption{Endpoints REST disponibles}
\begin{tabularx}{\linewidth}{llX}
  \toprule
  \textbf{M\'ethode} & \textbf{Endpoint} & \textbf{Description} \\
  \midrule
  GET & \texttt{/wp-json/tma/v1/destinations}       & Liste des destinations ACTIF \\
  GET & \texttt{/wp-json/tma/v1/offres}             & Offres publi\'ees avec m\'etadonn\'ees TMA \\
  GET & \texttt{/wp-json/tma/v1/client/\{id\}/fidelite} & Profil fid\'elit\'e d'un client \\
  \bottomrule
\end{tabularx}
\end{table}

\section{Multilingue --- Int\'egration Polylang}

\begin{phpbox}
// Enregistrement des chaines traduisibles
add_action( 'init', function() {
    if ( function_exists('pll_register_string') ) {
        pll_register_string( 'Reserver maintenant',
            'Reserver maintenant', 'Travel Maroc Theme' );
        // Options traduisibles (multilignes)
        $about = get_option('tma_about_text', '');
        if ( $about )
            pll_register_string('tma_about_text', $about,
                'Travel Maroc Options', true);
    }
});

// Lecture d'une option dans la langue courante
function tma_get_option_translated( string $key, string $default = '' ): string {
    $value = get_option( $key, $default );
    if ( function_exists('pll__') ) return pll__( $value ) ?: $value;
    return $value;
}
\end{phpbox}

Les m\'etadonn\'ees TMA num\'eriques sont \textbf{synchronis\'ees automatiquement}
entre les traductions via le filtre \texttt{pll\_copy\_post\_metas}~:
\texttt{\_tma\_points}, \texttt{\_tma\_guide\_id}, \texttt{\_tma\_hotel\_id},
\texttt{\_tma\_transport\_id}.

\section{Notifications Automatiques}

\begin{table}[H]
\centering
\caption{Notifications g\'en\'er\'ees automatiquement}
\begin{tabularx}{\linewidth}{llX}
  \toprule
  \textbf{\'Ev\'enement} & \textbf{Type} & \textbf{Canal} \\
  \midrule
  Nouvelle commande   & CONFIRMATION\_RESERVATION & EMAIL \\
  Commande termin\'ee & POINTS\_GAGNES             & EMAIL \\
  Mont\'ee de palier  & MONTEE\_PALIER             & EMAIL \\
  Commande annul\'ee  & ANNULATION                & EMAIL \\
  \bottomrule
\end{tabularx}
\end{table}

Les emails sont envoy\'es par lot via WP-Cron (fr\'equence~: toutes les heures,
hook~: \texttt{tmb\_send\_notifications\_cron}).

% ───────────────────────────────────────────────────────
\chapter{S\'ecurit\'e de l'Application}
% ───────────────────────────────────────────────────────

\section{Mod\`ele de S\'ecurit\'e WordPress}

L'application applique les cinq piliers de s\'ecurit\'e WordPress~:

\begin{enumerate}
  \item \textbf{Nonces} : Chaque formulaire et lien d'action inclut un nonce v\'erifi\'e
        c\^ot\'e serveur (\texttt{wp\_nonce\_field}, \texttt{check\_admin\_referer},
        \texttt{wp\_verify\_nonce})
  \item \textbf{Capabilities} : Chaque action v\'erifie la capacit\'e de l'utilisateur
        (\texttt{manage\_options}, \texttt{manage\_woocommerce}, capacit\'es TMA custom)
  \item \textbf{Sanitisation} : Toutes les entr\'ees utilisateur sont nettoy\'ees
        (\texttt{sanitize\_text\_field}, \texttt{sanitize\_email}, \texttt{esc\_url\_raw})
  \item \textbf{\'Echappement} : Toutes les sorties HTML utilisent \texttt{esc\_html()},
        \texttt{esc\_attr()}, \texttt{esc\_url()}, \texttt{wp\_kses\_post()}
  \item \textbf{Requ\^etes pr\'epar\'ees} : Toutes les requ\^etes SQL utilisent
        \texttt{\$wpdb->prepare()} pour pr\'evenir les injections SQL
\end{enumerate}

\section{Capacit\'es Personnalis\'ees}

\begin{phpbox}
// Classe TMB_Roles
const CAPS = [
    'tma_view_dashboard',   // Voir le dashboard TMA
    'tma_view_clients',     // Acceder a la liste clients
    'tma_view_orders',      // Acceder aux commandes
    'tma_manage_bookings',  // Gerer les reservations
];

public static function add_roles(): void {
    // Role Agent TMA (lecture seule)
    add_role( 'tma_agent', 'Agent TMA',
        array_merge(['read' => true],
            array_fill_keys(self::CAPS, true)) );
    // Ajouter les caps au Shop Manager
    $shop_mgr = get_role('shop_manager');
    foreach ( self::CAPS as $cap ) {
        $shop_mgr->add_cap( $cap );
    }
}
\end{phpbox}

\section{Protection contre les Injections SQL}

\begin{phpbox}
// Toutes les requetes parametrees utilisent $wpdb->prepare()
$profile = $wpdb->get_row( $wpdb->prepare(
    "SELECT p.*, t.taux_remise, t.points_jamais_expirent
     FROM {$wpdb->prefix}tma_client_profile p
     JOIN {$wpdb->prefix}tma_type_client t ON t.id = p.type_client_id
     WHERE p.wp_user_id = %d",
    $user_id
));
\end{phpbox}

\section{S\'ecurit\'e des Actions Destructives}

Lors de la \textbf{d\'esinstallation} du plugin, le hook \texttt{register\_uninstall\_hook}
d\'eclenche \texttt{TMB\_Activator::uninstall()} qui~:
\begin{itemize}
  \item Supprime les 10 tables TMA avec \texttt{DROP TABLE IF EXISTS}
  \item Supprime toutes les options \texttt{tma\_*}
  \item Annule les crons planifi\'es
\end{itemize}

La d\'esactivation (sans d\'esinstallation) ne supprime \textbf{aucune donn\'ee}.

% ───────────────────────────────────────────────────────
\chapter{Guide d'Installation et d'Utilisation}
% ───────────────────────────────────────────────────────

\section{Pr\'erequis}

\begin{itemize}
  \item \textbf{XAMPP} v8.2+ (Apache + MySQL + PHP 8.2)~: \url{https://www.apachefriends.org}
  \item \textbf{WordPress} 6.x install\'e dans \texttt{htdocs/boutique}
  \item \textbf{WooCommerce} 10.x activ\'e
  \item Navigateur moderne (Chrome recommand\'e)
\end{itemize}

\section{Installation en 5 \'Etapes}

\subsection{\'Etape 1 --- D\'emarrer XAMPP}

\begin{shellbox}
# Lancer XAMPP Control Panel
# Demarrer Apache + MySQL
# Verifier : http://localhost/boutique/ accessible
\end{shellbox}

\subsection{\'Etape 2 --- Activer le Plugin}

\begin{shellbox}
# WP Admin > Extensions > travel-maroc-booking > Activer
# L'activation cree automatiquement :
#   - 10 tables tma_* dans MySQL
#   - 5 paliers de fidelite (Standard, Silver, Gold, VIP, Entreprise)
#   - 16 localisations marocaines avec coordonnees GPS
#   - 4 guides, 6 hotels, 5 transports de demonstration
#   - Role tma_agent + capabilities shop_manager
#   - Crons : expiration points (quotidien) + notifications (horaire)
\end{shellbox}

\subsection{\'Etape 3 --- Configurer les R\'eglages TMA}

\begin{shellbox}
# WP Admin > Travel Maroc > Reglages
# Renseigner :
#   - URL video hero (.mp4)
#   - Image hero (fallback)
#   - Telephone, Email, WhatsApp, Adresse
#   - ID Google Analytics (G-XXXXXXXXXX)
#   - Reseaux sociaux (Facebook, Instagram)
#   - Texte page A-propos
\end{shellbox}

\subsection{\'Etape 4 --- Configurer Polylang (Multilingue)}

\begin{shellbox}
# WP Admin > Langues > Ajouter une langue
#   - Francais (fr) > definir comme defaut
#   - Arabe (ar) > cocher "Langue de droite a gauche"
# Produits > Selectionner tout > Action "Set language: Francais"
# Langues > Traductions de chaines > Traduire en arabe
\end{shellbox}

\subsection{\'Etape 5 --- Cr\'eer un compte Shop Manager}

\begin{shellbox}
# WP Admin > Travel Maroc > Reglages > Comptes & Roles
# Cliquer "Creer le compte Shop Manager"
# Mot de passe affiche une seule fois -- noter immediatement
# Login : shop_manager_tma
\end{shellbox}

\section{Points d'\'Evaluation}

\begin{table}[H]
\centering
\caption{Points d'\'evaluation sugg\'er\'es}
\begin{tabularx}{\linewidth}{lXc}
  \toprule
  \textbf{Fonctionnalit\'e} & \textbf{URL / Action} & \textbf{Statut} \\
  \midrule
  Page d'accueil + carte Leaflet   & \texttt{/boutique/}               & $\checkmark$ \\
  Catalogue + filtres cat\'egorie  & \texttt{/boutique/boutique-offres/} & $\checkmark$ \\
  Fiche offre + ressources         & \texttt{/boutique/product/...}    & $\checkmark$ \\
  Checkout + date/voyageurs        & \texttt{/boutique/commander/}     & $\checkmark$ \\
  Remise fid\'elit\'e automatique  & Panier client Silver/Gold/VIP     & $\checkmark$ \\
  Mon compte + r\'eservations      & \texttt{/boutique/mon-compte/}    & $\checkmark$ \\
  Notifications client             & Dashboard Mon compte              & $\checkmark$ \\
  Dashboard admin + KPIs           & TMA > Dashboard                   & $\checkmark$ \\
  Export CSV r\'eservations        & TMA > Dashboard > bouton          & $\checkmark$ \\
  CRUD Destinations                & TMA > Destinations                & $\checkmark$ \\
  CRUD Guides / H\^otels / Transport & TMA > Guides / H\^otels / Transport & $\checkmark$ \\
  Calendrier disponibilit\'e       & Produit > m\'eta box Disponibilit\'es & $\checkmark$ \\
  Sélecteur de langue              & Header (FR / AR)                  & $\checkmark$ \\
  Support RTL arabe                & Changer langue en arabe           & $\checkmark$ \\
  Open Graph + JSON-LD             & Voir source HTML d'une offre      & $\checkmark$ \\
  Sitemap XML                      & \texttt{/boutique/wp-sitemap.xml} & $\checkmark$ \\
  Journal d'audit                  & TMA > Journal d'audit             & $\checkmark$ \\
  \bottomrule
\end{tabularx}
\end{table}

\begin{infobox}
\textbf{Compte de d\'emonstration} \\
Admin WordPress~: identifiants configurés \`a l'installation.\\
Shop Manager~: \texttt{shop\_manager\_tma} / mot de passe g\'en\'er\'e
(visible une seule fois dans TMA > R\'eglages > Comptes).
\end{infobox}

% ───────────────────────────────────────────────────────
\chapter{Bilan et Perspectives}
% ───────────────────────────────────────────────────────

\section{Fonctionnalit\'es R\'ealis\'ees}

\begin{table}[H]
\centering
\caption{Bilan complet des fonctionnalit\'es}
\begin{tabularx}{\textwidth}{lXc}
  \toprule
  \textbf{Fonctionnalit\'e} & \textbf{D\'etails} & \textbf{Statut} \\
  \midrule
  Catalogue d'offres        & Filtres, cartes produit, m\'etadonn\'ees TMA     & $\checkmark$ \\
  Checkout enrichi          & Date, voyageurs, demandes sp\'eciales           & $\checkmark$ \\
  Programme de fid\'elit\'e & 4 paliers, remise auto, expiration cron         & $\checkmark$ \\
  Calendrier dispo          & Meta box admin, blocage dates, validation serveur & $\checkmark$ \\
  Espace client             & Dashboard, r\'eservations enrichies, notifs     & $\checkmark$ \\
  Dashboard admin           & 6 KPIs, top offres, logs, export CSV            & $\checkmark$ \\
  CRUD m\'etier             & Destinations, Guides, H\^otels, Transport       & $\checkmark$ \\
  Multilingue FR/AR         & Polylang, sélecteur header, RTL CSS             & $\checkmark$ \\
  Carte interactive         & Leaflet.js + OpenStreetMap                      & $\checkmark$ \\
  Animation logo            & GSAP 3 (draw + fade)                           & $\checkmark$ \\
  SEO complet               & OG, Twitter Cards, JSON-LD, Sitemap XML        & $\checkmark$ \\
  Google Analytics 4        & Snippet GA4 conditionnel (opt-in admin)        & $\checkmark$ \\
  S\'ecurit\'e              & Nonces, capabilities, sanitisation, audit log  & $\checkmark$ \\
  Emails personnalis\'es    & WC emails avec d\'etails voyage + WhatsApp     & $\checkmark$ \\
  API REST                  & 3 endpoints JSON (destinations, offres, fid.) & $\checkmark$ \\
  Uninstall propre          & Drop tables + delete options + clear crons     & $\checkmark$ \\
  \bottomrule
\end{tabularx}
\end{table}

\section{Difficult\'es Rencontr\'ees}

\begin{enumerate}
  \item \textbf{Polylang et produits sans langue} : Apr\`es activation de Polylang
        et ajout des langues, les produits cr\'e\'es ant\'erieurement n'avaient pas
        de langue assign\'ee et \'etaient masqu\'es. Solution~: action groupée
        dans la liste des produits pour assigner la langue.

  \item \textbf{Ordre du seed data} : Les tables \texttt{tma\_guide} et
        \texttt{tma\_hotel} r\'ef\'erencent \texttt{tma\_localisation} via cl\'e
        \'etrang\`ere. Le seed des guides/h\^otels devait imp\'erativement
        s'ex\'ecuter \textit{apr\`es} le seed des localisations.

  \item \textbf{dbDelta et requ\^etes multiples} : La fonction \texttt{dbDelta()}
        de WordPress ne g\`ere pas les requ\^etes multiples dans une seule cha\^ine
        SQL. Solution~: d\'ecoupage en requ\^etes individuelles via
        \texttt{explode(';', \$sql)}.

  \item \textbf{WooCommerce HPOS} : La v10.7 de WooCommerce utilise par d\'efaut
        le syst\`eme HPOS (\textit{High-Performance Order Storage}) qui stocke les
        commandes dans \texttt{wc\_orders} et non plus \texttt{wp\_posts}.
        Les requ\^etes SQL du dashboard ont \'et\'e adapt\'ees en cons\'equence.

  \item \textbf{IP spoofing dans le journal d'audit} : L'impl\'ementation initiale
        lisait \texttt{HTTP\_X\_FORWARDED\_FOR}, modifiable par n'importe quel
        client. Correction~: usage exclusif de \texttt{REMOTE\_ADDR}.
\end{enumerate}

\section{Perspectives d'Am\'elioration}

\begin{enumerate}
  \item \textbf{Passerelle de paiement} --- Int\'egration CMI (paiement en ligne
        Maroc) ou Stripe pour r\'eservations payables directement en ligne
  \item \textbf{Filtre avanc\'e boutique} --- Barre de recherche par destination,
        fourchette de prix et dur\'ee via \texttt{pre\_get\_posts}
  \item \textbf{Application mobile} --- PWA ou application React Native
        consommant l'API REST TMA
  \item \textbf{D\'eploiement cloud} --- Migration vers un h\'ebergeur en ligne
        (OVH, PlanetHoster, Hostinger) avec certificat SSL
  \item \textbf{Tableau de bord analytique} --- Graphiques Chart.js int\'egr\'es
        dans le dashboard admin (CA mensuel, top destinations)
  \item \textbf{Tests automatis\'es} --- Suite de tests PHP avec PHPUnit pour
        les classes du plugin
\end{enumerate}

% ───────────────────────────────────────────────────────
\chapter{Conclusion}
% ───────────────────────────────────────────────────────

Ce projet final a permis de concevoir et d\'evelopper \textbf{Travel Maroc Agency},
une plateforme e-commerce compl\`ete et op\'erationnelle pour une agence de voyage
marocaine. Le projet couvre l'int\'egralit\'e de la cha\^ine de valeur num\'erique~:
de la pr\'esentation publique multilingue des offres jusqu'\`a l'administration
m\'etier avanc\'ee, en passant par la r\'eservation, la fid\'elisation client et
la journalisation des actions.

\bigskip
Les principaux apports de ce projet sont~:

\begin{itemize}
  \item \textbf{Ma\^itrise de WordPress/WooCommerce} : Architecture hooks/filters,
        Plugin PHP orient\'e objet (8 classes), Template Hierarchy, overrides WC
  \item \textbf{Conception BDD relationnelle} : Mod\'elisation de 10 tables MySQL
        sp\'ecialis\'ees avec \texttt{dbDelta()}, cl\'es \'etrang\`eres et
        indexation
  \item \textbf{S\'ecurit\'e applicative} : Nonces, capabilities, sanitisation,
        \'echappement, requ\^etes pr\'epar\'ees, anti-IP spoofing
  \item \textbf{Int\'egration multilingue} : Polylang avec API native PHP, support
        RTL complet pour l'arabe, synchronisation des m\'etadonn\'ees
  \item \textbf{SEO et visibilit\'e} : Open Graph, Twitter Cards, Schema.org
        JSON-LD TouristTrip, Sitemap XML natif WordPress, GA4
  \item \textbf{UX et design} : Th\`eme sur mesure, animations GSAP, carte
        Leaflet.js, espace client enrichi avec tableau de bord fid\'elit\'e
\end{itemize}

\bigskip
Ce projet d\'emontre la capacit\'e \`a construire une application web professionnelle
et maintenable en utilisant l'\'ecosyst\`eme WordPress/WooCommerce, tout en
allant bien au-del\`a des fonctionnalit\'es par d\'efaut gr\^ace \`a un plugin
m\'etier PHP d\'evelopp\'e de z\'ero.

% ─────────────────────────────────────────────────────────────
%  BIBLIOGRAPHIE
% ─────────────────────────────────────────────────────────────
\chapter*{Bibliographie et R\'ef\'erences}
\addcontentsline{toc}{chapter}{Bibliographie et R\'ef\'erences}
\thispagestyle{fancy}

\section*{Documentation Officielle}
\begin{itemize}
  \item \textbf{WordPress Developer} --- \url{https://developer.wordpress.org}
  \item \textbf{WooCommerce Docs} --- \url{https://developer.woocommerce.com}
  \item \textbf{Polylang} --- \url{https://polylang.pro/doc}
  \item \textbf{Leaflet.js} --- \url{https://leafletjs.com/reference.html}
  \item \textbf{GSAP} --- \url{https://gsap.com/docs/v3}
  \item \textbf{Schema.org TouristTrip} --- \url{https://schema.org/TouristTrip}
  \item \textbf{Open Graph Protocol} --- \url{https://ogp.me}
  \item \textbf{Google Analytics 4} --- \url{https://developers.google.com/analytics}
\end{itemize}

\section*{Ressources}
\begin{itemize}
  \item MDN Web Docs --- APIs navigateur (CSS Variables, Fetch, localStorage)
  \item WordPress Codex --- Fonctions PHP WordPress (sanitisation, hooks, nonces)
  \item OWASP Top 10 --- Bonnes pratiques s\'ecurit\'e web
  \item OpenStreetMap --- Fournisseur de tuiles cartographiques open-source
\end{itemize}

% ─────────────────────────────────────────────────────────────
%  ANNEXES
% ─────────────────────────────────────────────────────────────
\appendix

\chapter{Structure Compl\`ete des Fichiers}

\begin{shellbox}
wp-content/
+-- plugins/travel-maroc-booking/
|   +-- travel-maroc-booking.php   # Point d'entree, register_*_hook
|   +-- includes/
|   |   +-- class-tmb-activator.php     # Tables, seed, crons, uninstall
|   |   +-- class-tmb-fidelite.php     # Points, paliers, remise panier
|   |   +-- class-tmb-checkout.php     # Champs voyage + validation
|   |   +-- class-tmb-availability.php # Meta box calendrier + AJAX
|   |   +-- class-tmb-notifications.php # File notifications + cron
|   |   +-- class-tmb-rest-api.php     # Endpoints /wp-json/tma/v1/*
|   |   +-- class-tmb-roles.php        # Roles + capabilities
|   |   +-- class-tmb-log-admin.php    # Journal d'audit
|   +-- admin/
|       +-- class-tmb-admin.php        # Dashboard + export CSV
|       +-- class-tmb-resources-admin.php # CRUD Guides/Hotels/Transport
|       +-- class-tmb-product-meta.php # Meta box produit (G/H/T + points)
|       +-- class-tmb-settings.php     # Page reglages TMA
+-- themes/travel-maroc/
    +-- style.css  functions.php  header.php  footer.php
    +-- index.php  page.php  single.php  404.php
    +-- page-about.php  page-contact.php
    +-- assets/css/main.css  logo-animation.css
    +-- assets/js/main.js  logo-animation.js
    +-- template-parts/
    |   +-- offer-card.php  about-content.php  contact-content.php
    +-- woocommerce/
    |   +-- single-product.php  archive-product.php
    |   +-- content-product.php  cart/cart.php
    |   +-- checkout/form-checkout.php  thankyou.php
    |   +-- myaccount/dashboard.php  orders.php  navigation.php
    |   +-- emails/customer-processing-order.php
    |   +-- emails/customer-completed-order.php
    +-- wpml-config.xml
    +-- languages/
\end{shellbox}

\chapter{Variables d'Environnement et Configuration}

\begin{shellbox}
# wp-config.php (extrait)
define( 'DB_NAME',     'boutique' );
define( 'DB_USER',     'root' );
define( 'DB_PASSWORD', '' );
define( 'DB_HOST',     'localhost' );
define( 'WP_DEBUG',    false );

# URL de base
define( 'WP_SITEURL', 'http://localhost/boutique' );
define( 'WP_HOME',    'http://localhost/boutique' );
\end{shellbox}

\chapter{Commandes Essentielles}

\begin{shellbox}
# == Verification syntaxe PHP ===========================
php -l wp-content/plugins/travel-maroc-booking/travel-maroc-booking.php
php -l wp-content/themes/travel-maroc/functions.php

# == Demarrage XAMPP ====================================
# Via XAMPP Control Panel : Start Apache + MySQL

# == Acces base de donnees ==============================
# http://localhost/phpmyadmin/
# Base : boutique | Tables TMA : tma_*

# == Reinitialiser les tables TMA =======================
# WP Admin > Extensions > Desactiver > Reactivite plugin TMB

# == Seed donnees demo (sans reactivation) ==============
# WP Admin > Travel Maroc > Reglages > "Inserer les donnees demo"

# == Test API REST ======================================
# GET http://localhost/boutique/wp-json/tma/v1/destinations
# GET http://localhost/boutique/wp-json/tma/v1/offres
\end{shellbox}

\chapter{Tables MySQL TMA --- DDL Complet}

\begin{shellbox}
-- Crees par dbDelta() a l'activation du plugin

CREATE TABLE wp_tma_localisation (
  id INT UNSIGNED NOT NULL AUTO_INCREMENT,
  nom VARCHAR(100) NOT NULL,
  type VARCHAR(20) NOT NULL DEFAULT 'VILLE',
  code_pays VARCHAR(3),
  latitude DECIMAL(9,6), longitude DECIMAL(9,6),
  actif TINYINT(1) NOT NULL DEFAULT 1,
  PRIMARY KEY (id)
);

CREATE TABLE wp_tma_client_profile (
  id INT UNSIGNED NOT NULL AUTO_INCREMENT,
  wp_user_id BIGINT UNSIGNED NOT NULL,
  type_client_id INT UNSIGNED NOT NULL DEFAULT 1,
  points_fidelite_solde INT UNSIGNED NOT NULL DEFAULT 0,
  points_cumules_total INT UNSIGNED NOT NULL DEFAULT 0,
  statut VARCHAR(10) NOT NULL DEFAULT 'ACTIF',
  date_inscription DATETIME NOT NULL DEFAULT CURRENT_TIMESTAMP,
  PRIMARY KEY (id),
  UNIQUE KEY uq_wp_user (wp_user_id)
);

-- + 8 autres tables : tma_destination, tma_guide, tma_hotel,
--   tma_transport, tma_type_client, tma_historique_points,
--   tma_notification, tma_log_admin
\end{shellbox}

% ─────────────────────────────────────────────────────────────
%  PIED DE PAGE FINAL
% ─────────────────────────────────────────────────────────────
\vspace{2cm}
\begin{center}
\begin{tikzpicture}
  \fill[estblue] (0,0) rectangle (\textwidth,1cm);
  \node[white,font=\small,align=center] at (0.5\textwidth,0.5cm)
    {Rapport r\'edig\'e par JOUMAIL Abderrahmane\\
     CDL 1\textsuperscript{\`ere} ann\'ee \textbullet{} EST Dakhla \textbullet{} 2025--2026};
\end{tikzpicture}
\end{center}

\end{document}
